\documentclass[10pt]{article}
\usepackage[a4paper,margin=1in]{geometry}
\usepackage[T1]{fontenc}
\usepackage{lmodern}
\usepackage{amsmath,amssymb,amsthm}
\usepackage{microtype}
\usepackage[colorlinks=true,linkcolor=blue,urlcolor=blue,citecolor=blue]{hyperref}
\newtheorem*{proposition}{Proposition}
\title{A Product-Form Counterexample Answers Q3 and Q5\\
\large Literature-status note for arXiv:1607.08024}
\author{}
\date{11 August 2026}
\begin{document}
\maketitle

\paragraph{Status.}
The two questions below have negative answers, implied by a later explicit
example of An--Lai \cite{AL}.  This note records the exact identification; it
does not claim a new result.

\section*{The source questions}
Dutkay--Haussermann--Lai \cite[Section 10, Q3 and Q5, PDF p.~38]{DHL} ask,
for an expansive integer matrix $R$ and integer digit set $B$:
\begin{quote}
\textbf{Q3.} If the self-affine measure $\mu(R,B)$ is spectral, must there be
$L$ such that $(R,B,L)$ is a Hadamard triple?

\textbf{Q5.} If $\mu(R,B)$ is spectral and has no overlap, must such an $L$
exist?
\end{quote}

\section*{The later example}
An--Lai \cite[Example 8.4, PDF p.~33]{AL} take
\[
   N=24,\qquad D=\{0,1,16,17\}.
\]
Their Theorem 2.6 proves that $\mu_{24,D}$ is a spectral measure.  Example
8.4 then observes that $D$ is not a spectral set in the finite cyclic group
$\mathbb Z_{24}$ and concludes explicitly that there is no integer set
$L$ for which $(24,D,L)$ is an ordinary Hadamard triple.

There is one extra point needed for Q5.  The elements of $D$ are pairwise
distinct modulo $24$, so $D$ is a simple digit set for $24$.  Theorem 1.6 of
the source paper \cite[PDF p.~5]{DHL} states that every simple digit set for
an expansive integer matrix satisfies the no-overlap condition.  Hence the
same example obeys the additional hypothesis in Q5.

\begin{proposition}
Questions Q3 and Q5 of \cite{DHL} both have negative answers.
\end{proposition}
\begin{proof}
For $(R,B)=(24,D)$, spectrality and failure of every ordinary Hadamard triple
are supplied by \cite[Theorem 2.6 and Example 8.4]{AL}.  Simplicity of $D$
modulo $24$ and \cite[Theorem 1.6]{DHL} give no overlap.  Thus the example
refutes Q3 and also Q5.
\end{proof}

\paragraph{Scope.}
This does not answer Q4, which asks whether spectrality by itself forces
no overlap, nor the broad Fourier-frame question Q2.

\begin{thebibliography}{9}
\bibitem{DHL}
D. Dutkay, J. Haussermann, and C.-K. Lai,
\emph{Hadamard triples generate self-affine spectral measures},
Trans. Amer. Math. Soc. 371 (2019), 1439--1481;
\href{https://arxiv.org/abs/1607.08024}{arXiv:1607.08024}.

\bibitem{AL}
L.-X. An and C.-K. Lai,
\emph{Product-form Hadamard triples and its spectral self-similar measures},
\href{https://arxiv.org/abs/2209.05616}{arXiv:2209.05616}.
\end{thebibliography}
\end{document}

